% ==============================================================================
% 00_frontmatter.tex – title page, abstract, lists
% (thesis data and the look of the title page are in preamble.tex)
% ==============================================================================

% --- Title page ----------------------------------------------------------------
\BHTTitelseite

% --- Front matter: Roman page numbers --------------------------------------------
\pagenumbering{Roman}

% --- Abstract ------------------------------------------------------------------
\chapter*{Abstract}
\addcontentsline{toc}{chapter}{Abstract}
Webcam-based eye tracking is inexpensive and available on almost every laptop.
With generic models, however, it is often inaccurate, because people, cameras
and environments differ strongly. This thesis investigates how a pretrained
vision model can be adapted to a single person with only a small amount of
calibration data.

A ResNet-18 pretrained on ImageNet serves as the feature extractor. It is small
enough for a CPU, and its general image features are well established; it is
therefore also a common backbone in gaze estimation. The backbone stays frozen,
and only a small regression head is trained. This keeps the number of trainable
parameters low, which reduces the risk of overfitting with few calibration
images. In addition, the features of each image are computed only once, so a
training run takes only seconds. The head is first trained on the GazeCapture
dataset (base model) and then adapted with own webcam recordings from 18
sessions of one person. This personalization is needed because eye shape, head
pose, camera and screen differ from person to person, and a generic model does
not compensate for these differences. A few minutes of personal data are
enough, instead of a large dataset of one's own. Accuracy is measured with the
diagonal error, i.e.\ the error in percent of the screen diagonal.

On the own recordings, the base model is not better than a constant prediction
of the screen centre (30.6\,\% versus 30.3\,\%); it therefore mainly serves as
a starting point. With a realistic, temporal split of the data (training on the
beginning, test on the end of a recording), personalization reduces the error
from 31.8\,\% to 27.9\,\%, i.e.\ by about 12\,\%. In sessions with good
recording conditions the gain is much larger (46 to 47\,\%) and also holds for
a new recording. With a random split of the frames the error drops to
19.8\,\% (about one third). This value is too optimistic, however, because
almost identical neighbouring frames end up in training and test. One training
run takes only a few seconds on a CPU. A short moving average makes the
predictions much calmer, but does not improve the accuracy. Background and
lighting have a larger influence than the design of the calibration. Removing wrongly labelled frames after target jumps
improves the result only slightly.

The results show that data-efficient personalization is possible with simple
means. Its benefit, however, must be assessed with a realistic data split. In
addition, two sessions of a second person were recorded. With the whole camera
image as input, personalization brought no gain for her, and a calibration of
the first person did not transfer to her. In the web demo, which uses an eye
crop, personalization did work for her as well (error \SI{9.7}{\percent}
instead of \SI{26.9}{\percent}). The input is therefore decisive for whether
personalization succeeds for a person.

\Schlagwoerter{Keywords}{eye tracking, gaze estimation, webcam,
personalization, transfer learning, ResNet-18, GazeCapture}

% --- German abstract (only with \ZweiteSprachetrue, see preamble.tex) -----------
\ifZweiteSprache
\begin{otherlanguage}{ngerman}
\chapter*{Kurzfassung}
\addcontentsline{toc}{chapter}{Kurzfassung}
Webcam-basiertes Eye-Tracking ist kostengünstig und auf fast jedem Laptop
verfügbar. Mit allgemeinen Modellen ist es aber oft ungenau, weil sich Personen,
Kameras und Umgebungen stark unterscheiden. Diese Arbeit untersucht, wie sich
ein vortrainiertes Bildverarbeitungsmodell mit wenigen Kalibrierungsdaten an
eine einzelne Person anpassen lässt.

Als Feature Extractor dient ein auf ImageNet vortrainiertes ResNet-18. Es ist
klein genug für eine CPU, und seine allgemeinen Bild-Features sind gut erprobt;
deshalb wird es auch in der Blickschätzung häufig als Backbone verwendet. Der
Backbone bleibt eingefroren, und trainiert wird nur ein kleiner Regression
Head. So gibt es nur wenige trainierbare Parameter, was die Gefahr von
Overfitting bei wenigen Kalibrierbildern verringert. Außerdem werden die
Features jedes Bildes nur einmal berechnet, und ein Training dauert nur
Sekunden. Der Head wird zuerst auf dem Datensatz GazeCapture trainiert
(Basismodell) und danach mit eigenen Webcam-Aufnahmen aus 18 Sitzungen einer
Person angepasst. Diese Personalisierung ist nötig, weil sich Augenform,
Kopfhaltung, Kamera und Bildschirm von Person zu Person unterscheiden und ein
allgemeines Modell diese Unterschiede nicht ausgleicht. Dafür genügen wenige
Minuten persönlicher Daten statt eines großen eigenen Datensatzes. Als
Fehlermaß dient der Diagonalfehler, also der Fehler in Prozent der
Bildschirmdiagonale.

Auf den eigenen Aufnahmen ist das Basismodell nicht besser als eine konstante
Schätzung der Bildschirmmitte (30,6\,\% gegenüber 30,3\,\%); es dient daher
vor allem als Startpunkt. Bei einer realistischen, zeitlichen Aufteilung der
Daten (Training auf dem Anfang, Test auf dem Ende einer Aufnahme) senkt die
Personalisierung den Fehler von 31,8\,\% auf 27,9\,\%, also um etwa 12\,\%.
In Sitzungen mit guten Aufnahmebedingungen ist der Gewinn deutlich größer (46
bis 47\,\%) und gilt auch für eine neue Aufnahme. Bei einer zufälligen
Aufteilung der Bilder sinkt der Fehler auf 19,8\,\% (etwa ein Drittel). Dieser
Wert ist aber zu optimistisch, weil fast gleiche Nachbarbilder in Training und
Test liegen. Ein Trainingslauf dauert auf einer CPU nur wenige Sekunden. Ein kurzer Moving
Average macht die Vorhersagen deutlich ruhiger, verbessert die Genauigkeit aber
nicht. Hintergrund und Beleuchtung haben einen größeren Einfluss als die
Gestaltung der Kalibrierung. Das Entfernen falsch gelabelter Bilder nach
Zielsprüngen verbessert das Ergebnis nur wenig.

Die Ergebnisse zeigen: Eine dateneffiziente Personalisierung ist mit einfachen
Mitteln möglich. Ihr Nutzen muss aber mit einer realistischen Aufteilung der
Daten bewertet werden. Alle Aufnahmen stammen von einer Person; die Ergebnisse
gelten daher zunächst nur für diese Person.

\Schlagwoerter{Schlagwörter}{Eye-Tracking, Blickschätzung, Webcam,
Personalisierung, Transfer Learning, ResNet-18, GazeCapture}
\end{otherlanguage}
\fi

% --- Lists ---------------------------------------------------------------------
\tableofcontents
\listoffigures
\listoftables

% --- Abbreviations and symbols -------------------------------------------------
\chapter*{List of Abbreviations}
\addcontentsline{toc}{chapter}{List of Abbreviations}
{\renewcommand{\arraystretch}{1.08}
\begin{longtable}{@{}p{2.4cm}p{\dimexpr\textwidth-2.4cm-2\tabcolsep\relax}@{}}
  \toprule
  \textbf{Abbreviation} & \textbf{Meaning} \\
  \midrule
  \endhead
  \bottomrule
  \endfoot
  AdamW   & Adam optimizer with decoupled weight decay \\
  CI      & Confidence interval \\
  CLIP    & Contrastive Language-Image Pre-training \\
  CNN     & Convolutional neural network \\
  CPU     & Central processing unit \\
  DA      & Domain adaptation \\
  DG      & Domain generalization \\
  DINOv2  & Self-Distillation with No Labels, version 2 \\
  FC      & Fully connected (layer) \\
  fps     & Frames per second \\
  L1      & Loss based on the mean absolute error \\
  LoRA    & Low-Rank Adaptation \\
  MAE     & Masked autoencoder \\
  MAML    & Model-Agnostic Meta-Learning \\
  MLP     & Multi-layer perceptron \\
  PEFT    & Parameter-efficient fine-tuning \\
  PoR     & Point of regard (gaze point on the screen) \\
  ReLU    & Rectified linear unit \\
  RMSE    & Root mean squared error \\
  SD      & Standard deviation \\
  SVR     & Support vector regression \\
  TTA     & Test-time adaptation \\
  ViT     & Vision Transformer \\
  \midrule
  \multicolumn{2}{@{}l@{}}{\textbf{Symbols}} \\[0.3ex]
  $\Ediag$ & Diagonal error: Euclidean RMSE in percent of the screen diagonal \\
\end{longtable}}

% --- Main part: Arabic page numbers --------------------------------------------
\cleardoublepage
\pagenumbering{arabic}
